\section{An explicit asymptotic for the actual symmetric velocity trace}

\begin{theorem}[NN-A013: every fixed polynomial residue moment]
Under A there is a sequence $T_n\in(n,n+1]$ avoiding all zero ordinates
such that, for every fixed integer $r\ge0$ and every $\epsilon>0$,
\[
 \mathcal V_r(T_n):=
 \sum_{0<\Im\rho<T_n}
 \operatorname{Res}_{s=\rho}
 \frac{s^r\zeta(s+1)}{\zeta(s)}
 =\frac{i^rT_n^{r+1}}{2\pi(r+1)}
 +O_{r,\epsilon}(T_n^{r+\epsilon}).
\]
The sum is over distinct nontrivial zeros, with the full residue at a
multiple zero. No simplicity assumption is needed. If all these zeros
are simple, then
\[
 \mathcal V_r(T_n)=\sum_{0<\gamma<T_n}\rho^{r-1}v_\rho,
 \qquad v_\rho=\frac{\rho\zeta(\rho+1)}{\zeta'(\rho)}.
\]
In particular,
\[
 \mathcal V_0(T_n)=\frac{T_n}{2\pi}+O_\epsilon(T_n^\epsilon),
 \qquad
 \mathcal V_1(T_n)=\frac{iT_n^2}{4\pi}
 +O_\epsilon(T_n^{1+\epsilon}).
\]
The latter is the derivative of the sum of zeros in the fixed contour
under Hurwitz deformation, including ramified clusters as symmetric
traces. It is not an asymptotic for individual velocities at multiple
zeros.
\end{theorem}
\begin{proof}
We explicitly import the following standard consequence of RH: there
exist $T_n\in(n,n+1]$ and an absolute $A_0$ such that
\[
 \sup_{1/2\le\sigma\le2}|\zeta(\sigma+iT_n)|^{-1}
 \ll\exp(A_0\log T_n/\log\log T_n).
\]
A primary paper recording the precise statement is Milinovich--Ng,
\emph{A note on a conjecture of Gonek}, Section 3, equation (3.2),
printed p.5; its original locator is Titchmarsh, Theorem 14.16.
This imported bound assumes RH, not simplicity. The functional equation
and Stirling's formula extend it to $-1/2\le\sigma\le2$ as
\[
 |\zeta(\sigma+iT_n)|^{-1}
 \ll T_n^{\min(\sigma-1/2,0)+\epsilon}
\]
for every fixed $\epsilon>0$. A also gives the usual Lindelof bound
$\zeta(\sigma+it)\ll_\epsilon t^\epsilon$ uniformly on fixed compact
real intervals with $\sigma\ge1/2$, away from bounded $t$.

Integrate $f_r(s)=s^r\zeta(s+1)/\zeta(s)$ counterclockwise around
the rectangle with vertical sides $\Re s=2$ and $\Re s=-1/2$,
bottom at a fixed positive nonzero-ordinate height $T_0$, and top at
$T_n$. Its only poles are the nontrivial zeta zeros inside; the pole
of $\zeta(s+1)$ is at $s=0$, outside the rectangle, and $1/\zeta(s)$
has a zero rather than a pole at $s=1$. Trivial zeros are outside.
Adding the fixed finitely many omitted lower residues contributes $O(1)$.

On the right edge,
\[
 \frac{\zeta(s+1)}{\zeta(s)}=
 \sum_{m\ge1}b(m)m^{-s},\qquad
 b(m)=m^{-1}\prod_{p\mid m}(1-p),\quad b(1)=1.
\]
The bound $|b(m)|\le1$ gives absolute convergence. The $m=1$ term
contributes
\[
 \frac1{2\pi}\int_{T_0}^{T_n}(2+it)^r\,dt
 =\frac{i^rT_n^{r+1}}{2\pi(r+1)}+O_r(1+T_n^r).
\]
For $m\ge2$, integration by parts yields
$|\int_{T_0}^{T_n}(2+it)^re^{-it\log m}dt|
\le C_r(1+T_n)^r/\log m$.
Summing with $|b(m)|m^{-2}$ gives an $O_r(1+T_n^r)$ error.

On the left edge, the functional equation and absolute convergence at
$3/2-it$ give $|1/\zeta(-1/2+it)|\ll t^{-1}$.
The numerator is $\zeta(1/2+it)\ll t^\epsilon$, so
$|f_r(-1/2+it)|\ll t^{r-1+\epsilon}$.
Its integral is $O_{r,\epsilon}(T_n^{r+\epsilon})$, including $r=0$.
On the top edge the numerator lies in $\Re(s+1)\ge1/2$, hence is
$O(T_n^\epsilon)$, and the imported good-height reciprocal estimate
gives $f_r=O(T_n^{r+2\epsilon})$ uniformly. The horizontal length
is fixed. The bottom edge is a constant contribution. Replace
$\epsilon$ by $\epsilon/2$ at the start and apply the residue theorem.
At a simple zero the residue equals
$\rho^r\zeta(\rho+1)/\zeta'(\rho)=\rho^{r-1}v_\rho$.
At a multiple zero the exact interpretation is NN-A005.
\end{proof}

These formulas strengthen the elementary contour identity by supplying a
verified asymptotic at explicit good-height sequences, and repair the
missing simplicity qualification in a velocity interpretation. They do
not extend automatically to every real cutoff $T$: an individual residue
near a cutoff need not be bounded by the displayed total error. Nor does
the small real part of $\mathcal V_1$ imply equal proportions of positive
and negative real velocities. A sum can cancel with arbitrarily unequal
sign counts; $N-1$ entries equal to $1$ and one equal to $-(N-1)$ already
show this. The odd phase-correlation gate of NN-A002 remains a separate
arithmetic problem.
